In this paper we shall solve the generalized problem for complex finite-dimensional representations of finite groups.
Let $G$ be a finite group.
Let $\rh\colon G \to \on{GL}(V)$ be a representation of~$G$ on~a~finite-dimensional complex vector space $V$. By~Hilbert's finiteness theorem the algebra of~inva\-riant polynomials $\C[V]^G$ is finitely generated.
Let $\si_1,\dots,\si_n$ be a system of generators, we~call them \emph{basic invariants}, and let
$\si= (\si_1,\dots,\si_n)$ be the resulting map $\si \colon V \to \C^n$.
The map~$\si$ separates $G$-orbits and hence induces a homeomorphism between the orbit space $V/G$ and the image $\si(V)$.
(Notice that since $G$ is finite and thus all $G$-orbits are closed, there is a bijection between the orbits and
the points in the affine variety $V \cq G$ with coordinate ring $\C[V]^G$; in other words the \emph{categorical quotient}
$V \cq G$ is a \emph{geometric quotient}.)
As a consequence we may identify $V/G$ with $\si(V)$ and the canonical orbit projection $V \to V/G$ with $\si \colon V \to \si(V)$. We~will also write $G \acts V$ for the representation $\rh$.


The basic invariants can be chosen to be homogeneous polynomials. A system of homogeneous basic invariants is \emph{minimal} if
none among them is superfluous. In~that case their number and their degrees are uniquely determined (cf.\ \cite[p.~95]{DK02}).

Assume that a map $f \colon \Om \to \si(V)$ defined on some open subset $\Om \subseteq \R^m$ is given. We~will assume that $f$ possesses some degree of differentiability as a map into $\C^n$.
The question we will address in this paper is the following:
\begin{quote}
 How \emph{differentiable} can lifts of $f$ over $\si$ be? By a \emph{lift of $f$ over $\si$} we mean a map
 $\ol f \colon \Om \to V$ such that $f = \si \o \ol f$.
\end{quote}
Simple examples show that, in general, a big loss of regularity occurs from $f$ to lifts of $f$. We~will determine the optimal regularity of lifts among the Sobolev spaces $W^{1,p}$ under minimal differentiability
requirements on $f$. In~particular, the optimal $p>1$ will be determined as an~explicit function of the maximal homogeneity degree
of the basic invariants.

Note that the results do not depend on the choice of the basic invariants since any two choices differ by a polynomial
diffeomorphism.
